\documentclass{ermacrs2023}

\title{XI Encontro Regional de Matemática Aplicada e Computacional do Rio Grande do Sul -- Instruções para Trabalhos Completos}

% Usar \author{}, com o argumento em branco, para suprimir os nomes
% \author{}
\author{Nome Sobrenome$^{1}$, \underline{Nome Sobrenome}$^{2}$, Nome Sobrenome$^{3}$, Nome Sobrenome$^{4}$}

% Usar \info{}, com o argumento em branco, para suprimir as filiações e os endereços eletrônicos
% \info{}
\info{$^1$Universidade, xxxx@yyy.zzz; \\ $^2$Universidade, aaaa@bbb.ccc; \\$^3$Universidade, nnnn@ppp.qqq; \\ $^4$Universidade, nnnn@ppp.qqq.}
  
\abstract{Os \textbf{trabalhos completos} submetidos para apreciação da Comissão Científica do XI ERMAC--RS devem seguir o padrão apresentado neste documento modelo. O resumo deve conter, no máximo, 200 palavras. Ele  deve enunciar de forma clara e sintética: problema de pesquisa, abordagem metodológica empreendida, resultados e conclusões.}
  
\keywords{Palavra 1; Palavra 2; Palavra 3; Palavra 4.}

\begin{document}
  \maketitle
  
  \section{Introdução}\label{sec:introducao}
  Os \textbf{trabalhos completos} devem ter de 5 a 8 páginas, ser redigidos em \textbf{língua portuguesa} e estar de acordo com este modelo modelo \LaTeX.
  
  A submissão de trabalhos deverá ocorrer pela plataforma \texttt{Even3}, em endereço eletrônico divulgado na página do evento. O trabalho deve ser submetido em formato \texttt{.pdf}. Também será solicitada a submissão de uma versão do trabalho não identificada. Não serão aceitas propostas enviadas por qualquer outro meio, tampouco após os prazos.

  \subsection{Sessões Técnicas}\label{subsec:sessoes-tecnicas}
  
  No momento da submissão os trabalhos devem ser classificados em uma das sessões técnicas:
  \begin{itemize}
    \item ST1 – Análise Aplicada
    \item ST2 – Biomatemática
    \item ST3 – Física-Matemática, Mecânica dos Fluidos e Sistemas Dinâmicos
    \item ST4 – Matemática Aplicada à Engenharia
    \item ST5 – Modelagem Matemática e Computacional
    \item ST6 – Métodos Estocásticos e Estatísticos
    \item ST7 – Métodos Numéricos e Computação Científica
    \item ST8 – Matemática Discreta e Otimização
    \item ST9 – Ensino e Extensão em Matemática Aplicada e Computacional
    \item ST10 – Análise
    \item ST11 – Geometria
    \item ST12 – Álgebra
    \item ST13 – Teoria dos Números
    \item ST14 – Modelagem Matemática Aplicada às Ciências Naturais
  \end{itemize}

  \subsection{Submissão de Trabalhos}\label{subsec:submissao}

  A data final para a submissão dos trabalhos é \textbf{31/03/2023}, com divulgação dos resultados ocorrendo até o dia \textbf{30/04/2023}.
  
  Demais informações sobre o evento e sobre a submissão de trabalhos podem ser encontradas no site \href{https://wp.ufpel.edu.br/ermacrs23}{\textbf{\underline{https://wp.ufpel.edu.br/ermacrs23}}}. 
  
  \section{Seções do Trabalho}\label{sec:secoes-trabalho}
  
  Sugere-se que o trabalho seja organizado nas seguintes seções: Introdução, metodologia, resultados e discussões, conclusões, agradecimentos (se houver) e referências. As seções ``AGRADECIMENTOS'' e ``REFERÊNCIAS'' não devem ser numeradas.
  
  \section{Como inserir  equações}\label{sec:equacoes} 
  
  Todas as equações apresentadas no texto devem ser sequencialmente numeradas, conforme os exemplos a seguir
  \begin{equation}\label{eq:pvi}
    \begin{cases}
      y^{\prime}(t) = f(t,\,y(t)),\\
      y(t_{0}) = y_{0},
    \end{cases} 
  \end{equation}
  e
  \begin{equation}\label{eq:equacao-quadratica}
    ax^{2} + bx + c = 0.
  \end{equation}
 
  \section{Figuras e Tabelas}\label{sec:figuras-tabelas}

  As figuras e as tabelas apresentadas no corpo do texto devem estar próximas do local onde são inicialmente referidas.

  \subsection{Figuras}\label{subsec:figuras}
  
  Para inserir figuras no corpo do texto utilize o ambiente \texttt{figure}. Todas as figuras presente no texto deverem ser numeradas, centralizadas e ter a sua fonte adequadamente referenciada. Ainda, elementos textuais presentes nas figuras devem ser legíveis, com tamanho de fonte compatível à do corpo do texto. Por fim, recomenda-se que as figuras possuam uma resolução mínima de 300 dpi. Como exemplo, temos a Figura~\ref{fig:logo-ermacrs} a seguir.
  \begin{figure}[ht]
    \centering
    \includegraphics[width=0.5\textwidth]{logoermacrs.jpeg}
    \caption{Exemplo de figura. \\ Fonte: do autor.} 
    \label{fig:logo-ermacrs}
  \end{figure}

  \subsection{Tabelas}\label{subsec:tabelas}

  Para inserir tabelas no corpo do texto utilize o ambiente \texttt{table}. Todas as tabelas presentes no texto devem ser numeradas, centralizadas e ter a sua fonte adequadamente referenciada. Ainda, a legenda da tabela deve estar acima dela. A Tabela~\ref{tab:exemplo-tabela} é um exemplo de tabela e segue o padrão \LaTeX.

  \begin{table}[ht]
   \centering
   \caption{Exemplo de tabela$^{*}$.} \label{tab:exemplo-tabela}
   \begin{tabular}{ccc}
      \toprule
        Elemento X & Elemento Y & Elemento Z \\ \midrule
        0,1 & 0,3 & 0,4 \\ \bottomrule
    \end{tabular}
    \caption*{Fonte: do autor. \\ $^{*}$Alguma informação adicional sobre a tabela, se tiver.}
  \end{table}

  \section{Seção de Referências}

  Na seção de referências, devem estar listados todos as obras citadas, por ordem alfabética de sobrenomes. Para citações, os autores devem utilizar comandos como o \texttt{citep} e o \texttt{citeyear}, conforme os exemplos:: \citep{strogatz2015,EadesWormald,Brusquetta}, ``ainda Viana (\citeyear{Viana}) diz que...''.

  \textbf{Esta seção tem um limite de tamanho de uma página. Além disso, esta seção não contabiliza no número total de páginas de cada artigo.}
  
  \section{Conclusões}
  
  O trabalho deve conter um tópico de conclusão apresentando o fechamento do trabalho.
  
  \section*{Agradecimentos}
  
  Agradecimentos, se houver, devem vir logo após a última seção do corpo do texto e antes da seção “REFERÊNCIAS”.
  
  \begin{thebibliography}{99}
    \bibitem[BRUSQUETTA {\it et al.}, 2019]{Brusquetta}Brusqueta, M., Picotti E., Mion, E., Chen, Y., Beghi, A., Minen, D. A., ``Nonlinear Model Predictive Control based Virtual Driver for high performance driving''. Em: {\bf Annals of IEEE Conference on Control Technology and Applications} (2019), Hong Kong, China.
      
    \bibitem[EADES; WORMALD,  1994]{EadesWormald} Eades, P.,  Wormald, N. C.  ``Edge crossings in drawings of bipartite graphs''. Em: {\bf Algorithmica} 11 (1994), pp. 379--403. DOI: \texttt{10.1007/BF01187020}.

    \bibitem[XI ERMARC-RS, 2023]{Ermac} ERMAC--RS. Site oficial do XI Congresso Regional de Matemática Aplicada e Computacional do Rio Grande do Sul. Online. Acessado em 07/02/2023, \href{https://wp.ufpel.edu.br/ermacrs23}{\textbf{\underline{https://wp.ufpel.edu.br/ermacrs23}}}.
      
    \bibitem[STROGATZ,  2015]{strogatz2015}Strogatz, S. H. {\bf Nonlinear Dynamics and Chaos: With Applications to Physics, Biology, Chemistry, and Engineering},  2nd ed. Philadelphia: Westview Press, 2015. ISBN: \texttt{978-0813349107}.
      
    \bibitem[VIANA, 2019]{Viana} Viana, D. S.  ``Condições de otimalidade para otimização cônica''. Tese de doutorado. Instituto de Matemática e estatística, Universidade de São Paulo, Programa de Matemática Aplicada, 2019.
  \end{thebibliography}
\end{document}
